On considère une solution aqueuse $\mathrm{(S_{1})}$ d'acide méthanoïque
$\ce{HCOOH}$ de concentration molaire ${C_{1}=\qty{5e-3}{\mol\per\L}}$. La
mesure de la conductivité de cette solution, à $\qty{25}{\degreeCelsius}$, donne
$\sigma_{1}=\qty{33}{\mS\per\m}$.

\begin{donnees}
    \begin{itemize}
        \item Conductivités molaires ioniques~:

              $\lambda_{1}=\lambda(\ce{H3O^+})=\qty{35.0}{\mS\square\m\per\mol}$~;
              $\lambda_{2}=\lambda(\ce{HCOO^-})=\qty{5.5}{\mS\square\m\per\mol}$

        \item L'effet des ions $\ce{HO^{-}_{(aq)}}$ sur la conductivité de la
              solution est supposé négligeable.
        \item La conductivité $\sigma$ d'une solution s'écrit en fonction des
              concentrations molaires effectives des ions $\ce{X_{i}}$ et de
              leurs conductivités molaires ioniques $\lambda_{i}$ comme suit~:
              $\sigma=\sum\lambda_{i}\bigl[\ce{X_{i}}\bigr]$.
    \end{itemize}
\end{donnees}
\begin{enumerate}
    \item Écrire l'équation chimique modélisant la réaction entre l'acide
          méthanoïque et l'eau.
    \item Montrer que
          $\bigl[\ce{H3O^{+}_{(aq)}}\bigr]=\qty{8.15e-4}{\mol\per\L}$.
    \item Calculer le taux d'avancement final $\tau_{1}$ de la réaction.
          Conclure.
    \item Montrer que la valeur du quotient de réaction à l'état d'équilibre du
          système chimique est ${Q_{r_{1},\eq}=\num{1.59e-4}}$.
    \item On dilue à $\qty{25}{\degreeCelsius}$ la solution $\mathrm{(S_{1})}$,
          pour obtenir une solution aqueuse $\mathrm{(S_{2})}$ de concentration
          molaire $C_{2}$.
    
          Donner, en justifiant, la valeur du quotient de réaction
          $Q_{r_{2},\eq}$ du système chimique à l'état d'équilibre.
\end{enumerate}


